
\documentclass[a4paper, 10pt]{article}
\usepackage{a4wide}
\usepackage {amsmath}
\usepackage {graphicx}
\usepackage[utf8]{inputenc} 
\usepackage[francais]{babel}
\usepackage{fancyhdr}
\usepackage{setspace}
\usepackage{fancyhdr}
\usepackage{lastpage}
\usepackage{extramarks}
\usepackage{chngpage}
\usepackage{soul}
\usepackage[usenames,dvipsnames]{color}
\usepackage{graphicx,float,wrapfig}
\usepackage{ifthen}
\usepackage{listings}
\usepackage{courier}
\usepackage{stmaryrd}
\usepackage[T1]{fontenc}
\usepackage{mathrsfs}

\lhead{} 
\chead{} 
\rhead{\bfseries Intégration Numérique} 
\lfoot{JC Toussaint}
%\cfoot{} 
%\rfoot{\thepage}

% This is the color used for MATLAB comments below
\definecolor{MyDarkGreen}{rgb}{0.0,0.4,0.0}

% For faster processing, load Matlab syntax for listings
\lstloadlanguages{Matlab}%
\lstset{language=Matlab,                        % Use MATLAB
        frame=single,                           % Single frame around code
        basicstyle=\small\ttfamily,             % Use small true type font
        keywordstyle=[1]\color{Blue}\bf,        % MATLAB functions bold and blue
        keywordstyle=[2]\color{Purple},         % MATLAB function arguments purple
        keywordstyle=[3]\color{Blue}\underbar,  % User functions underlined and blue
        identifierstyle=,                       % Nothing special about identifiers
                                                % Comments small dark green courier
        commentstyle=\usefont{T1}{pcr}{m}{sl}\color{MyDarkGreen}\small,
        stringstyle=\color{Purple},             % Strings are purple
        showstringspaces=false,                 % Don't put marks in string spaces
        tabsize=5,                              % 5 spaces per tab
        %
        %%% Put standard MATLAB functions not included in the default
        %%% language here
        morekeywords={xlim,ylim,var,alpha,factorial,poissrnd,normpdf,normcdf},
        %
        %%% Put MATLAB function parameters here
        morekeywords=[2]{on, off, interp},
        %
        %%% Put user defined functions here
        morekeywords=[3]{FindESS, homework_example},
        %
        morecomment=[l][\color{Blue}]{...},     % Line continuation (...) like blue comment
        numbers=left,                           % Line numbers on left
        firstnumber=1,                          % Line numbers start with line 1
        numberstyle=\tiny\color{Blue},          % Line numbers are blue
        stepnumber=5                            % Line numbers go in steps of 5
        }

% Includes a MATLAB script.
% The first parameter is the label, which also is the name of the script
%   without the .m.
% The second parameter is the optional caption.
\newcommand{\matlabscript}[2]
  {\begin{itemize}\item[]\lstinputlisting[caption=#2,label=#1]{#1.m}\end{itemize}}

\pagestyle{fancy}

\begin{document}

\title{Méthode de quadrature de Gauss-Legendre}

%\author{Jean-Christophe Toussaint\\
%  Phelma\\
 %\texttt{jean-christophe.toussaint@grenoble-inp.fr}}
\date{\today}
 
\maketitle

\section{Polynomes de Legendre}

Le polynome de Legendre de degré $n$ est solution de l'équation différentielle suivante:

\begin{equation}
\frac{d}{du} \big[(1-u^2) \frac{d}{du} P_n(u) \big]+n(n+1) P_n(u)=0
\label{eqdif_legendre}
\end{equation}
où $u \in [-1, 1]$. Il vérifie $P_n(1)=1$.

On peut générer la suite des polynomes de Legendre en utilisant la formule de récurrence de Bonnet
pour tout entier $n>0$:
\begin{equation}
(n+1) \; P_{n+1}(u) = (2n+1)\; u \; P_n(u) -n \; P_{n-1}(u)
\label{recurrence}
\end{equation}
sachant que $P_0(u)=1$ et $P_1(u)=u$.

Les  polynomes de Legendre forment une base orthogonale :
\begin{equation}
\langle P_m, \; P_n \rangle=\int_{-1}^{1} P_m(u) \; P_n(u)\; du = \frac{2}{2n+1}\delta_{m,n}
\label{ortho}
\end{equation}

\begin{enumerate}
\item Utiliser le programme Matlab {\tt main\_legendre} pour tracer $P_n(u)$ pour  $n  \llbracket  0, 3  \rrbracket$.
\end{enumerate}

\section{Polynomes de Lagrange}

On veut construire un polynome de degré $n$ passant par $n+1$ points de coordonnées $(x_i, y_i)$ avec $i \in  \llbracket  0, n  \rrbracket $,
aussi appelé interpolée. Son expression est donnée par :
\begin{equation}
\mathscr{I}_n(x)= \sum_{i=0}^n y_i \; L_i(x)
\label{interpolee}
\end{equation}

où $L_i(x) = \prod_{j=0, j \ne i}^n \,\frac{ (x-x_j) }{ (x_i-x_j) }$ est le polynome de Lagrange de degré $n$ associé au point $i$.

En imposant $\mathscr{I}_n(x_i) \equiv y_i$  en tout point d'échantillonnage $i$, la propriété suivante est obligatoirement vérifiée:
\begin{equation}
L_i(x_j)=\delta_{i,j}
\label{unite}
\end{equation}

\textbf{Notez que les polynomes de Lagrange  forment une base mais non orthogonale.}

\section{Quadrature de Gauss-Legendre}

On échantillonne la fonction $f(x)$ aux abscisses $x_i$ où $i \in  \llbracket  0, n  \rrbracket $.
On peut l'écrire sous la forme du développement suivant :

\begin{equation}
f(x) = \mathscr{I}_n(x) + R_n(x) = \sum_{i=0}^{n} L_i(x) \, f(x_i) + \left[ \prod_{i=0}^n (x-x_i) \right] \frac{f^{(n+1)}(\xi(x))}{(n+1)!}
\label{interpol1}
\end{equation}
avec $\xi(x) \in [a, b]$
où $L_i(x)$ est le polynome de Lagrange associé au noeud $i$ et $R_n(x)$ est l'erreur due à l'approximation polynomiale de $f(x)$.

Le but ici est d'estimer numériquement l'intégrale $\int_a^b f(x) \; dx$ en approximant la fonction $f(x)$ par son interpolée $\mathscr{I}_n(x)$
de degré $n$. Formellement, on peut écrire:

\begin{equation}
I=\int_a^b\, f(x)\, dx=\int_a^b\, \mathscr{I}_n(x)\, dx + \int_a^b\, R_n(x)\, dx
\end{equation}

Sans perdre en généralité, l'intervalle d'intégration $[a, b]$ peut être ramené à $[-1, 1]$ par le changement de variable suivant :
$x = \frac{b-a}{2} \; u + \frac{a+b}{2}$ avec $u \in [-1, 1]$. On note dans la suite $g(u) = f(x(u))$.

L'expression \ref{interpol1} devient alors:

\begin{equation}
g(u) =  \sum_{i=0}^{n} L_i(u) \, g(u_i) + \left[ \prod_{i=0}^n (u-u_i) \right] \frac{g^{(n+1)}(\xi(u))}{(n+1)!}
\label{interpol2}
\end{equation}
avec $\xi(u) \in [-1, 1]$ où $u_i = \frac{2}{b-a}\;  (x_i-\frac{a+b}{2})$.

Si on suppose maintenant que $g(u)$ est un polynome de degré $2n+1$ alors le terme $\frac{g^{(n+1)}(\xi(u))}{(n+1)!}$ sera
de degré $n$ au plus, puisque $ \sum_{i=0}^{n} L_i(u) \, g(u_i)$ est un polynome de degré $n$ au plus et $\Pi_{i=0}^n (u-u_i)$ 
est de degré $n+1$.

Pour simplifier la lecture, on note $q_n(u) = \frac{g^{(n+1)}(\xi(u))}{(n+1)!}$ qui est un polynome de degré $n$. Alors
\begin{equation}
g(u) =  \sum_{i=0}^{n} L_i(u) \, g(u_i) + \left[ \prod_{i=0}^n (u-u_i) \right] q_n(u)
\label{interpol3}
\end{equation}

\subsection{Détermination des poids et des abscisses de Gauss}

Nous allons montrer qu'un choix judicieux des points d'échantillonnage permet d'annuler l'erreur d'intégration d'un polynome de degré $n$,
représenté par le dernier terme de l'équation précédente.

Pour ce faire, on intégre les deux cotés de l'équation \ref{interpol3} sur l'intervalle $[-1, 1]$, ce qui nous donne :

\begin{equation}
\int_{-1}^{1}\; g(u)\; du =  \int_{-1}^{1}\;  \sum_{i=0}^{n} L_i(u) \, g(u_i) \; du+ \int_{-1}^{1}\;  \left[ \prod_{i=0}^n (u-u_i) \right] q_n(u)\; du
\label{interpol4}
\end{equation}

On remarque $g(u_i)$ est une constante qui peut être sortie de l'intégrale. On obtient finalement

\begin{equation}
\int_{-1}^{1}\; g(u)\; du =   \sum_{i=0}^{n} w_i\, g(u_i)  + R
\label{interpol5}
\end{equation}
avec $w_i =  \int_{-1}^{1}\;  L_i(u) \; du$ est le poids de gauss associé au noeud $i$.
Dans cette relation,  le premier terme de droite est appelée formule de quadrature et le dernier le reste $R$.
Celui-ci s'écrit:

\begin{equation}
R = \int_{-1}^{1}\;  \left[ \prod_{i=0}^n (u-u_i) \right] q_n(u)\; du
\label{reste}
\end{equation}

Intéressons-nous maintenant au reste.
On développe dans un premier temps les deux polynomes $ \prod_{i=0}^n (u-u_i)$ et $q_n(u)$ sur la base des polynomes
de Legendre. On écrit formellement :

\begin{equation}
\prod_{i=0}^n (u-u_i)= \sum_{i=0}^{n+1} b_i \; P_i(u) \text{\quad et \quad} q_n(u) = \sum_{i=0}^{n} c_i \; P_i(u)
\label{expandR1}
\end{equation}
où $(b_i, \; c_i)$ sont des coefficients que l'on déterminera ci-après.

L'expression développée du dernier terme de l'équation \ref{interpol5} s'écrit alors:

\begin{equation}
\int_{-1}^{1}\; \left[ \sum_{i=0}^{n}  \sum_{j=0}^{n} \; b_i\; c_j P_i(u)\; P_j(u) + \sum_{i=0}^{n}\; b_{n+1} \;c_i P_i(u) \; P_{n+1}(u) \right] \; du
\label{expandR2}
\end{equation}

Grâce à la propriété d'orthogonalité des polynomes de Legendre, tous les termes de la forme $\; b_i\; c_j \int_{-1}^{1}\; P_i(u)\; P_j(u)\; du $
disparaissent si $i \ne j$.

L'équation \ref{expandR2} se simplifie en :
\begin{equation}
\int_{-1}^{1}\; \sum_{i=0}^{n}  \; b_i\; c_i \left[P_i(u) \right]^2\; \; du
\label{expandR3}
\end{equation}


Sachant que les $n$ abscisses $u_i$ ne sont pas fixées a priori,
une stratégie pour faire disparaitre ce terme d'erreur est  d'imposer que le produit suivant ne s'exprime
qu'en fonction de $P_{n+1}(u) $ : 
\begin{equation}
\prod_{i=0}^n (u-u_i)= b_{n+1} \; P_{n+1}(u) .
\label{prod1}
\end{equation}
Ce qui revient à faire coincider les abscisses $u_i$  pour  
$i \in  \llbracket  0, n  \rrbracket $ avec les racines du polynome de legendre  $P_{n+1}(u) $.
Comme les polynomes de legendre forment une base,  on en déduit que $b_i \equiv 0$ pour  
$i \in  \llbracket  0, n  \rrbracket $.

En égalisant les coefficients du monome $u^{n+1}$ avec celui du polynome  $P_{n+1}(u)$, on en déduit
la valeur de $b_{n+1}$.
Par exemple pour $n=3$, le coefficient du monome de plus haut degré de $P_4(u)$ est $\frac{35}{8}$, comme
celui de $\prod_{i=0}^n (u-u_i)$ est $1$ alors on en déduit que $b_{n+1}=\frac{8}{35}$.

On arrive finalement au dénouement de cette démonstration. En faisant coïncider les abscisses $u_i$ avec les zéros du 
polynome de Legendre $P_{n+1}(u)$, le terme d'erreur de la formule de quadrature
disparait complétement  lorsqu'on intégre une fonction polynomiale de degré $2n+1$.

On obtient finalement la relation de quadrature suivante :
\begin{equation}
\int_{-1}^{1}\; g(u)\; du =   \sum_{i=0}^{n} w_i\, g(u_i) 
\label{quadrature}
\end{equation}

Une analyse plus poussée [Hildebrand] permet d'établir une expression analytique pour les poids de Gauss associé au noeud $i$:
\begin{equation}
w_i =  \int_{-1}^{1}\;  L_i(u) \; du = \frac{2 (1-u_i^2)}{n\; [P_{n-1}(u_i)]^2}
\label{poids}
\end{equation}

% Une référence \cite{Press:2007}
 %\cite{Dhatt:2005}

%\bibliography{biblio}

\vspace{-0.5cm}
\section{Références}

\noindent Hildebrand F. B. : {\it Methods of Applied Mathematics }, ed. Dover Publications Inc., 1965.

\newpage

\section{Annexe}

Nous allons donner les éléments pour retrouver l'expression \ref{interpol1} utilisée dans le calcul de l'erreur d'interpolation $R_n$ d'une fonction par un polynome de degré $n$.
Comme $f(x) - P_n(x)-R_n(x)=0$ par définition, on a par conséquent :

\begin{equation}
 f(x) - P_n(x) - \left[ \prod_{i=0}^n (x-x_i) \right] q(x) =0
\label{err1}
\end{equation}

Dans le but de déterminer $q(x)$, on construit une fonction auxiliaire qu'on note $W(t)$ définie par :
\begin{equation}
W(t) =  f(t) - P_n(t) - \left[ \prod_{i=0}^n (t-x_i) \right] q(x) =0
\label{aux1}
\end{equation}
On remarque en particulier que $q(x)$ n'a pas été remplacé  par $q(t)$. La fonction $W(t)$ dépend en fait de $t$ et de $x$ mais
on ne s'intéresse ici qu'aux variations liées à $t$, la valeur de $x$ étant fixée.

On étudie maintenant les racines de $W(t)$. Il est évident que $W(t)$ s'annule en $t=x_0, x_1, ... , x_n$ mais aussi en $t=x$ par construction
d'après \ref{err1}.

 $W(t)$ doit être continue et dérivable pour appliquer le théorème de Rolle. Dans ce cas,  $W'(t)$ possède une racine entre chacune des
 $n+2$ zéros de $W(t)$, donc en tout $n+1$ zéros. Si $W''(t)$ existe ce que nous allons supposer, elle présente $n$ racines et ainsi de suite
 avec les dérivées successives de $W(t)$. Finalement $W^{(n+1)}(t)$ possède au moins une racine dans l'intervalle $[x_0, x_n]$ que l'on 
 note $t=\xi$.
 On obtient finalement l'identité suivante :
 
\begin{equation}
W^{(n+1)}(\xi)= 0 = \frac{d^{n+1}}{dt^{n+1}}  \left[ f(t) - P_n(t) - \left[ \prod_{i=0}^n (t-x_i) \right] q(x) \right]_{t=\xi}
\label{aux2}
\end{equation}

d'où en développant la dérivée d'ordre $n+1$, 

\begin{equation}
q(x) = \frac{f^{(n+1)}(\xi)}{(n+1)! }   .
\label{aux3}
\end{equation}


\end{document}

